% ===============================================================
%  Задачи 11-13 (объёмы: цилиндрические, сферические, обобщённые сферические координаты)
% ===============================================================

% ---------------------------------------------------------------
\task{Найти объём тела, ограниченного поверхностями $az=x^2+y^2$, $z=\sqrt{x^2+y^2}$ $(a>0)$.}

\step{1}{Шапка}
\begin{shapka}
$\displaystyle V=\iiint\limits_{V}dx\,dy\,dz$
\tcblower
$\partial V:\quad
\begin{array}{ll}
az=x^2+y^2 & (1)\\
z=\sqrt{x^2+y^2} & (2)
\end{array}\qquad a>0$
\end{shapka}

\step{2}{Два графика}
(1) --- параболоид вращения $z=\frac{x^2+y^2}{a}$ (чаша с вершиной в $O$), (2) --- верхняя половина конуса с углом $45^\circ$ к оси $Oz$. Обе поверхности --- поверхности вращения вокруг $Oz$, поэтому берём \textbf{цилиндрические координаты}: (1) $z=\dfrac{r^2}{a}$, (2) $z=r$.

Пересечение: $\dfrac{r^2}{a}=r\Rightarrow r=0$ или $r=a$ (окружность радиуса $a$ на высоте $z=a$). При $0<r<a$ имеем $\dfrac{r^2}{a}<r$: параболоид ниже конуса. Значит, (1) --- пол, (2) --- крыша, проекция --- круг $r\leqslant a$, а угол $\varphi$ пробегает полный оборот. Тело --- «линза» между конусом и чашей.

\begin{center}
\begin{minipage}[t]{0.49\linewidth}\centering
\figEleven

\small График 1: конус (2) и параболоид (1) (вырезана четверть, чтобы видеть тело между ними); $a=1$
\end{minipage}\hfill
\begin{minipage}[t]{0.49\linewidth}\centering
\begin{tikzpicture}[scale=3.4,>=Stealth]
  \fill[newreg] (0,0)--(1,1) -- plot[domain=1:0,samples=30] (\x,{\x*\x}) -- cycle;
  \draw[->] (-0.08,0)--(1.3,0) node[below]{$r$};
  \draw[->] (0,-0.08)--(0,1.3) node[left]{$z$};
  \draw[very thick,newline] (0,0)--(1.15,1.15);
  \draw[very thick,newline] plot[domain=0:1.07,samples=30] (\x,{\x*\x});
  \node[above left] at (0.3,0.3) {\scriptsize (2) $z=r$};
  \node[below right] at (0.78,0.6) {\scriptsize (1) $z=\frac{r^2}{a}$};
  \draw[dashed,gray] (1,0)--(1,1)--(0,1);
  \node[below] at (1,0) {\scriptsize $a$};
  \node[left] at (0,1) {\scriptsize $a$};
  \draw[->,thick,hint] (0.6,0.36)--(0.6,0.6);
  \draw[dashed,hint] (0.6,0)--(0.6,0.36);
  \node[below] at (0.6,0) {\scriptsize $r$};
  \node at (0.42,0.3) {\scriptsize $\Omega_{rz}$};
\end{tikzpicture}

\small График 2: сечение в плоскости $(r,z)$; угол $0\leqslant\varphi\leqslant2\pi$ --- «прямоугольная» переменная
\end{minipage}
\end{center}

\step{3}{Замена переменных и якобиан}
\[
\begin{cases}x=r\cos\varphi,\\ y=r\sin\varphi,\\ z=z,\end{cases}\qquad |J|=r .
\]

\step{4}{Пределы по графику 2}
$0\leqslant\varphi\leqslant2\pi$;\quad $0\leqslant r\leqslant a$;\quad $\dfrac{r^2}{a}\leqslant z\leqslant r$ (вертикаль входит через параболоид, выходит через конус).

\step{5}{Повторный интеграл и счёт}
\[
V=\int_0^{2\pi}d\varphi\int_0^a r\,dr\int_{r^2/a}^{r}dz=2\pi\int_0^a\Bigl(r^2-\frac{r^3}{a}\Bigr)dr=2\pi\Bigl(\frac{a^3}{3}-\frac{a^3}{4}\Bigr)=\frac{2\pi a^3}{12}=\frac{\pi a^3}{6}.
\]

\begin{answerbox}
V=\frac{\pi a^3}{6}
\end{answerbox}

\emph{Геометрическая проверка.} Тело $=$ (чаша параболоида до высоты $a$) $-$ (конус высоты $a$). Объём чаши: сечение $z=\mathrm{const}$ --- круг радиуса $\sqrt{az}$, $\int_0^a\pi az\,dz=\frac{\pi a^3}{2}$. Объём конуса с радиусом основания $a$ и высотой $a$: $\frac13\pi a^2\cdot a=\frac{\pi a^3}{3}$. Разность: $\frac{\pi a^3}{2}-\frac{\pi a^3}{3}=\frac{\pi a^3}{6}$. \checkmark

% ---------------------------------------------------------------
\task{Переходя к сферическим или цилиндрическим координатам, вычислить объём, ограниченный поверхностями $x^2+y^2+z^2=2az$, $x^2+y^2\leqslant z^2$.}

\step{1}{Шапка}
\begin{shapka}
$\displaystyle V=\iiint\limits_{V}dx\,dy\,dz$
\tcblower
$\partial V:\quad
\begin{array}{ll}
x^2+y^2+z^2=2az & (1)\\
x^2+y^2=z^2 & (2)
\end{array}$\\[1mm]
$x^2+y^2\leqslant z^2$ (внутри конуса), $a>0$
\end{shapka}

\step{2}{Два графика}
(1): выделим полный квадрат: $x^2+y^2+(z-a)^2=a^2$ --- сфера радиуса $a$ с центром $(0,0,a)$; она касается плоскости $Oxy$ в начале координат и лежит в области $z\geqslant0$. (2) --- конус с углом $45^\circ$; условие $x^2+y^2\leqslant z^2$ --- «внутри конуса», и внутри шара это его верхняя половина.
Пересечение: подставим $x^2+y^2=z^2$ в (1): $2z^2=2az\Rightarrow z=a$ --- окружность радиуса $a$ на высоте $a$ (это «экватор» сферы). Тело --- «мороженое в рожке»: снизу конус от $O$ до высоты $a$, сверху --- полусфера.

Тело симметрично относительно оси $Oz$, поэтому достаточно нарисовать сечение плоскостью $y=0$. Удобнее всего \textbf{сферические координаты}: (1) $r^2=2ar\sin\psi\Rightarrow r=2a\sin\psi$; (2) $r^2\cos^2\psi\leqslant r^2\sin^2\psi\iff\operatorname{tg}\psi\geqslant1\iff\psi\geqslant\frac\pi4$.

\begin{center}
\begin{minipage}[t]{0.49\linewidth}\centering
\begin{tikzpicture}[scale=1.45,>=Stealth]
  \fill[reg] (0,0)--(1,1) arc[start angle=0,end angle=180,radius=1] -- cycle;
  \draw[->] (-1.6,0)--(1.7,0) node[below]{$x$};
  \draw[->] (0,-0.2)--(0,2.5) node[left]{$z$};
  \draw[thick,regline,dashed] (0,1) circle(1);
  \draw[very thick,regline] (1,1) arc[start angle=0,end angle=180,radius=1];
  \draw[very thick,regline] (-1,1)--(0,0)--(1,1);
  \draw[dashed,gray] (0,0)--(1.6,1.6) node[right,black]{\scriptsize (2)};
  \draw[dashed,gray] (0,0)--(-1.6,1.6) node[left,black]{\scriptsize (2)};
  \node[right] at (0.75,1.85) {\scriptsize (1)};
  \draw[dashed,gray] (-1,1)--(1,1);
  \node[right] at (1,0.92) {\scriptsize $z=a$};
  \node[left] at (0,2.08) {\scriptsize $2a$};
  \draw[->,very thick,hint] (0,0)--(0.766,1.643);
  \draw[hint] (0.45,0) arc(0:65:0.45);
  \node[hint] at (0.6,0.32) {\scriptsize $\psi$};
  \draw (0.3,0.3) arc(45:90:0.4243);
  \node at (0.2,0.62) {\tiny $\frac\pi4$};
  \node at (-0.45,1.2) {$V$};
\end{tikzpicture}

\small График 1: сечение плоскостью $y=0$ --- «рожок» (2) с «шапкой» (1)
\end{minipage}\hfill
\begin{minipage}[t]{0.49\linewidth}\centering
\begin{tikzpicture}
\begin{axis}[width=6.4cm,height=6.4cm,axis lines=left,xmin=0,xmax=1.8,ymin=0,ymax=2.4,
  xtick={0.7854,1.1345,1.5708},xticklabels={$\frac\pi4$,$\psi$,$\frac\pi2$},
  ytick={1.4142,2},yticklabels={$\sqrt2a$,$2a$},
  xlabel={$\psi$},ylabel={$r$},ylabel style={rotate=-90},clip=false]
  \addplot[name path=top,draw=none,domain=0.7854:1.5708,samples=40]{2*sin(deg(x))};
  \addplot[name path=bot,draw=none,domain=0.7854:1.5708]{0};
  \addplot[newreg] fill between[of=top and bot];
  \addplot[very thick,newline,domain=0.7854:1.5708,samples=40]{2*sin(deg(x))};
  \draw[very thick,newline] (axis cs:0.7854,0)--(axis cs:0.7854,1.4142);
  \draw[very thick,newline] (axis cs:1.5708,0)--(axis cs:1.5708,2);
  \draw[dashed,gray] (axis cs:0,1.4142)--(axis cs:0.7854,1.4142);
  \draw[dashed,gray] (axis cs:0,2)--(axis cs:1.5708,2);
  \draw[->,thick,hint] (axis cs:1.1345,0)--(axis cs:1.1345,1.8126);
  \node[anchor=east] at (axis cs:0.95,2.32) {\scriptsize (1) $r=2a\sin\psi$};
  \draw[->,gray] (axis cs:0.95,2.32)--(axis cs:1.0,1.72);
  \node[left] at (axis cs:0.7854,0.7) {\scriptsize (2)};
  \node at (axis cs:1.35,0.6) {\scriptsize $\Omega_{\psi r}$};
\end{axis}
\end{tikzpicture}

\small График 2: в плоскости $(\psi,r)$; долгота $0\leqslant\varphi\leqslant2\pi$
\end{minipage}
\end{center}

\step{3}{Замена переменных и якобиан}
\[
\begin{cases}x=r\cos\psi\cos\varphi,\\ y=r\cos\psi\sin\varphi,\\ z=r\sin\psi,\end{cases}\qquad |J|=r^2\cos\psi .
\]

\step{4}{Пределы по графику 2}
Сначала углы: $0\leqslant\varphi\leqslant2\pi$,\quad $\dfrac\pi4\leqslant\psi\leqslant\dfrac\pi2$; затем радиус: луч из $O$ под широтой $\psi$ идёт до сферы, $0\leqslant r\leqslant2a\sin\psi$.

\step{5}{Повторный интеграл и счёт}
\[
\begin{aligned}V&=\int_0^{2\pi}d\varphi\int_{\pi/4}^{\pi/2}\cos\psi\,d\psi\int_0^{2a\sin\psi}r^2\,dr
=2\pi\int_{\pi/4}^{\pi/2}\frac{8a^3\sin^3\psi}{3}\cos\psi\,d\psi\\
&=\frac{16\pi a^3}{3}\Bigl[\frac{\sin^4\psi}{4}\Bigr]_{\pi/4}^{\pi/2}
=\frac{16\pi a^3}{3}\cdot\frac14\Bigl(1-\frac14\Bigr)=\pi a^3 .\end{aligned}
\]
(Здесь $\sin\frac\pi4=\frac{\sqrt2}{2}$, $\sin^4\frac\pi4=\frac14$.)

\begin{answerbox}
V=\pi a^3
\end{answerbox}

\emph{Проверка 1 (геометрия).} Рожок --- конус с радиусом основания $a$ и высотой $a$: $\frac{\pi a^3}{3}$; шапка --- полушар радиуса $a$: $\frac23\pi a^3$. Сумма $\pi a^3$. \checkmark

\emph{Проверка 2 (цилиндрические координаты).} Сфера: $z=a\pm\sqrt{a^2-r^2}$; при $0\leqslant r\leqslant a$ тело занимает $r\leqslant z\leqslant a+\sqrt{a^2-r^2}$:
$V=2\pi\int_0^ar\bigl(a+\sqrt{a^2-r^2}-r\bigr)dr=2\pi\Bigl(\frac{a^3}{2}+\frac{a^3}{3}-\frac{a^3}{3}\Bigr)=\pi a^3$. \checkmark

% ---------------------------------------------------------------
\task{Пользуясь подходящей заменой переменных, вычислить объёмы тел, ограниченных поверхностями (параметры положительны):
а) $\sqrt{\dfrac xa}+\sqrt{\dfrac yb}+\sqrt{\dfrac zc}=1$ $(x,y,z\geqslant0)$;\quad
б) $\sqrt[3]{\dfrac xa}+\sqrt[3]{\dfrac yb}+\sqrt[3]{\dfrac zc}=1$ $(x,y,z\geqslant0)$;\quad
в) $\Bigl(\dfrac xa\Bigr)^{2/3}+\Bigl(\dfrac yb\Bigr)^{2/3}+\Bigl(\dfrac zc\Bigr)^{2/3}=1$.}

\textbf{Общая идея для всех трёх пунктов.} Все три поверхности имеют вид $\left(\frac xa\right)^{2/n}+\left(\frac yb\right)^{2/n}+\left(\frac zc\right)^{2/n}=1$: в а) $n=4$, в б) $n=6$, в в) $n=3$. Для них работают \textbf{обобщённые сферические координаты} (см.\ теорию):
\[
x=a\,r\cos^n\varphi\cos^n\psi,\quad y=b\,r\sin^n\varphi\cos^n\psi,\quad z=c\,r\sin^n\psi,
\]
\[
|J|=n^2abc\,r^2\cos^{n-1}\varphi\sin^{n-1}\varphi\cos^{2n-1}\psi\sin^{n-1}\psi,
\]
поверхность превращается в $r=1$, координатные плоскости --- в $\varphi=0$, $\varphi=\frac\pi2$, $\psi=0$, и часть тела в первом октанте становится параллелепипедом $[0,\frac\pi2]\times[0,\frac\pi2]\times[0,1]$.

\subsection*{а) $\sqrt{x/a}+\sqrt{y/b}+\sqrt{z/c}=1$}

\step{1}{Шапка}
\begin{shapka}
$\displaystyle V=\iiint\limits_{V}dx\,dy\,dz$
\tcblower
$\partial V:\quad
\begin{array}{ll}
\sqrt{\frac xa}+\sqrt{\frac yb}+\sqrt{\frac zc}=1 & (1)\\
x=0,\ y=0,\ z=0 & \text{(2)--(4)}
\end{array}$\\[1mm]
$x,y,z\geqslant0;\ a,b,c>0$
\end{shapka}

\step{2}{Два графика}
Поверхность (1) пересекает оси в точках $(a,0,0)$, $(0,b,0)$, $(0,0,c)$. Её след на плоскости $z=0$: $\sqrt{x/a}+\sqrt{y/b}=1$ --- дуга параболы, вогнутая к началу координат. Тело --- «криволинейный тетраэдр» с вогнутой гранью.

\begin{center}
\begin{minipage}[t]{0.49\linewidth}\centering
\figThirteen{4}

\small График 1: тело а) в первом октанте
\end{minipage}\hfill
\begin{minipage}[t]{0.49\linewidth}\centering
\BoxPic{\varphi}{\psi}{r}{\frac\pi2}{\frac\pi2}{1}

\small График 2: параллелепипед в $(\varphi,\psi,r)$
\end{minipage}
\end{center}

\step{3}{Замена переменных и якобиан ($n=4$)}
\[
\begin{cases}x=a\,r\cos^4\varphi\cos^4\psi,\\ y=b\,r\sin^4\varphi\cos^4\psi,\\ z=c\,r\sin^4\psi,\end{cases}
\qquad |J|=16\,abc\,r^2\cos^3\varphi\sin^3\varphi\cos^7\psi\sin^3\psi .
\]
Проверка: $\sqrt{\frac xa}+\sqrt{\frac yb}+\sqrt{\frac zc}=\sqrt r\,(\cos^2\varphi\cos^2\psi+\sin^2\varphi\cos^2\psi+\sin^2\psi)=\sqrt r$, т.е. (1) $\iff r=1$.

\step{4}{Пределы по графику 2}
$0\leqslant\varphi\leqslant\frac\pi2$,\quad $0\leqslant\psi\leqslant\frac\pi2$,\quad $0\leqslant r\leqslant1$.

\step{5}{Повторный интеграл и счёт}
\[
V=16abc\int_0^{\pi/2}\cos^3\varphi\sin^3\varphi\,d\varphi\cdot\int_0^{\pi/2}\cos^7\psi\sin^3\psi\,d\psi\cdot\int_0^1r^2dr .
\]
\begin{itemize}
\item $\displaystyle\int_0^{\pi/2}\sin^3\varphi\cos^3\varphi\,d\varphi\overset{t=\sin\varphi}{=}\int_0^1t^3(1-t^2)\,dt=\frac14-\frac16=\frac1{12}$ \ (так как $\cos^2\varphi=1-t^2$, $\cos\varphi\,d\varphi=dt$);
\item $\displaystyle\int_0^{\pi/2}\cos^7\psi\sin^3\psi\,d\psi\overset{t=\cos\psi}{=}\int_0^1t^7(1-t^2)\,dt=\frac18-\frac1{10}=\frac1{40}$ \ (так как $\sin^2\psi=1-t^2$, $\sin\psi\,d\psi=-dt$);
\item $\displaystyle\int_0^1r^2dr=\frac13$.
\end{itemize}
\[
V=16abc\cdot\frac1{12}\cdot\frac1{40}\cdot\frac13=\frac{16abc}{1440}=\frac{abc}{90}.
\]

\begin{answerbox}
V=\frac{abc}{90}
\end{answerbox}

\emph{Проверка другой заменой.} $x=aX^2$, $y=bY^2$, $z=cZ^2$ (треугольная замена, $|J|=2aX\cdot2bY\cdot2cZ=8abc\,XYZ$) переводит тело в тетраэдр $X+Y+Z\leqslant1$, $X,Y,Z\geqslant0$:
\[
\begin{aligned}V&=8abc\int_0^1X\,dX\int_0^{1-X}Y\,dY\int_0^{1-X-Y}Z\,dZ=8abc\int_0^1X\,dX\int_0^{1-X}\frac{Y(1-X-Y)^2}{2}dY\\
&=8abc\int_0^1\frac{X(1-X)^4}{24}dX=\frac{8abc}{24\cdot30}=\frac{abc}{90}.\end{aligned}
\]
(Использовано $\int_0^sY(s-Y)^2dY=\frac{s^4}{12}$ и $\int_0^1X(1-X)^4dX=\int_0^1(1-t)t^4dt=\frac15-\frac16=\frac1{30}$.) \checkmark

\clearpage
\subsection*{б) $\sqrt[3]{x/a}+\sqrt[3]{y/b}+\sqrt[3]{z/c}=1$}

\step{1}{Шапка}
\begin{shapka}
$\displaystyle V=\iiint\limits_{V}dx\,dy\,dz$
\tcblower
$\partial V:\quad
\begin{array}{ll}
\sqrt[3]{\frac xa}+\sqrt[3]{\frac yb}+\sqrt[3]{\frac zc}=1 & (1)\\
x=0,\ y=0,\ z=0 & \text{(2)--(4)}
\end{array}$\\[1mm]
$x,y,z\geqslant0;\ a,b,c>0$
\end{shapka}

\step{2}{Два графика}
Та же картина, что в а), только грань ещё сильнее «вдавлена» к началу координат (кубический корень растёт быстрее квадратного около нуля).

\begin{center}
\begin{minipage}[t]{0.49\linewidth}\centering
\figThirteen{6}

\small График 1: тело б) в первом октанте
\end{minipage}\hfill
\begin{minipage}[t]{0.49\linewidth}\centering
\BoxPic{\varphi}{\psi}{r}{\frac\pi2}{\frac\pi2}{1}

\small График 2: параллелепипед в $(\varphi,\psi,r)$
\end{minipage}
\end{center}

\step{3}{Замена переменных и якобиан ($n=6$)}
\[
\begin{cases}x=a\,r\cos^6\varphi\cos^6\psi,\\ y=b\,r\sin^6\varphi\cos^6\psi,\\ z=c\,r\sin^6\psi,\end{cases}
\qquad |J|=36\,abc\,r^2\cos^5\varphi\sin^5\varphi\cos^{11}\psi\sin^5\psi .
\]
Проверка: $\sqrt[3]{\frac xa}+\sqrt[3]{\frac yb}+\sqrt[3]{\frac zc}=\sqrt[3]r\,(\cos^2\varphi\cos^2\psi+\sin^2\varphi\cos^2\psi+\sin^2\psi)=\sqrt[3]r$, т.е. (1) $\iff r=1$.

\step{4}{Пределы по графику 2}
$0\leqslant\varphi\leqslant\frac\pi2$,\quad $0\leqslant\psi\leqslant\frac\pi2$,\quad $0\leqslant r\leqslant1$.

\step{5}{Повторный интеграл и счёт}
\[
V=36abc\int_0^{\pi/2}\cos^5\varphi\sin^5\varphi\,d\varphi\cdot\int_0^{\pi/2}\cos^{11}\psi\sin^5\psi\,d\psi\cdot\int_0^1r^2dr .
\]
\begin{itemize}
\item $\displaystyle\int_0^{\pi/2}\sin^5\varphi\cos^5\varphi\,d\varphi\overset{t=\sin\varphi}{=}\int_0^1t^5(1-t^2)^2dt=\int_0^1(t^5-2t^7+t^9)\,dt=\frac16-\frac14+\frac1{10}=\frac1{60}$;
\item $\displaystyle\int_0^{\pi/2}\cos^{11}\psi\sin^5\psi\,d\psi\overset{t=\cos\psi}{=}\int_0^1t^{11}(1-t^2)^2dt=\frac1{12}-\frac2{14}+\frac1{16}=\frac{28-48+21}{336}=\frac1{336}$;
\item $\displaystyle\int_0^1r^2dr=\frac13$.
\end{itemize}
\[
V=36abc\cdot\frac1{60}\cdot\frac1{336}\cdot\frac13=\frac{36abc}{60480}=\frac{abc}{1680}.
\]

\begin{answerbox}
V=\frac{abc}{1680}
\end{answerbox}

\clearpage
\subsection*{в) $(x/a)^{2/3}+(y/b)^{2/3}+(z/c)^{2/3}=1$}

\step{1}{Шапка}
\begin{shapka}
$\displaystyle V=\iiint\limits_{V}dx\,dy\,dz$
\tcblower
$\partial V:\quad
\Bigl(\frac xa\Bigr)^{2/3}+\Bigl(\frac yb\Bigr)^{2/3}+\Bigl(\frac zc\Bigr)^{2/3}=1\quad(1)$\\[1mm]
знаки $x,y,z$ любые; $a,b,c>0$
\end{shapka}

\step{2}{Два графика}
Здесь ограничений на знаки нет. Так как $\left(\frac xa\right)^{2/3}=\sqrt[3]{(x/a)^2}$ не меняется при $x\to-x$ (и так же для $y$, $z$), тело симметрично относительно всех трёх координатных плоскостей. Это «астроидальное» тело (трёхмерный аналог астроиды) с вершинами $(\pm a,0,0)$, $(0,\pm b,0)$, $(0,0,\pm c)$. По свойству симметрии
\[
V=8V_1,\qquad V_1\ \text{--- объём части в первом октанте.}
\]

\begin{center}
\begin{minipage}[t]{0.49\linewidth}\centering
\figAstroid

\small График 1: всё тело в); оно состоит из 8 одинаковых частей
\end{minipage}\hfill
\begin{minipage}[t]{0.49\linewidth}\centering
\BoxPic{\varphi}{\psi}{r}{\frac\pi2}{\frac\pi2}{1}

\small График 2: часть первого октанта в $(\varphi,\psi,r)$
\end{minipage}
\end{center}

\step{3}{Замена переменных и якобиан ($n=3$)}
\[
\begin{cases}x=a\,r\cos^3\varphi\cos^3\psi,\\ y=b\,r\sin^3\varphi\cos^3\psi,\\ z=c\,r\sin^3\psi,\end{cases}
\qquad |J|=9\,abc\,r^2\cos^2\varphi\sin^2\varphi\cos^5\psi\sin^2\psi .
\]
Проверка: $\left(\frac xa\right)^{2/3}+\left(\frac yb\right)^{2/3}+\left(\frac zc\right)^{2/3}=r^{2/3}(\cos^2\varphi\cos^2\psi+\sin^2\varphi\cos^2\psi+\sin^2\psi)=r^{2/3}$, т.е. (1) $\iff r=1$.

\step{4}{Пределы по графику 2 (первый октант)}
$0\leqslant\varphi\leqslant\frac\pi2$,\quad $0\leqslant\psi\leqslant\frac\pi2$,\quad $0\leqslant r\leqslant1$.

\step{5}{Повторный интеграл и счёт}
\[
V_1=9abc\int_0^{\pi/2}\cos^2\varphi\sin^2\varphi\,d\varphi\cdot\int_0^{\pi/2}\cos^5\psi\sin^2\psi\,d\psi\cdot\int_0^1r^2dr .
\]
\begin{itemize}
\item обе степени чётные --- понижаем степень:
\[\int_0^{\pi/2}\sin^2\varphi\cos^2\varphi\,d\varphi=\frac14\int_0^{\pi/2}\sin^22\varphi\,d\varphi=\frac14\int_0^{\pi/2}\frac{1-\cos4\varphi}{2}d\varphi=\frac{\pi}{16};\]
\item $\displaystyle\int_0^{\pi/2}\cos^5\psi\sin^2\psi\,d\psi\overset{t=\sin\psi}{=}\int_0^1(1-t^2)^2t^2dt=\int_0^1(t^2-2t^4+t^6)\,dt=\frac13-\frac25+\frac17=\frac{8}{105}$;
\item $\displaystyle\int_0^1r^2dr=\frac13$.
\end{itemize}
\[
V_1=9abc\cdot\frac{\pi}{16}\cdot\frac{8}{105}\cdot\frac13=\frac{\pi abc}{70},\qquad V=8V_1=\frac{8\pi abc}{70}=\frac{4\pi abc}{35}.
\]

\begin{answerbox}
V=\frac{4\pi abc}{35}
\end{answerbox}

\emph{Проверка на разумность:} эллипсоид с теми же полуосями имеет объём $\frac43\pi abc\approx4{,}19\,abc$, а наше тело с «вдавленными» гранями лежит внутри него: $\frac{4\pi}{35}abc\approx0{,}36\,abc$ --- меньше, как и должно быть. \checkmark

% ===============================================================
\clearpage
\section*{Сводка ответов}
\addcontentsline{toc}{section}{Сводка ответов}
\begin{center}
\renewcommand{\arraystretch}{2.5}
\begin{tabular}{>{\centering\arraybackslash}m{1.0cm} >{\raggedright\arraybackslash}m{7.4cm} >{\centering\arraybackslash}m{6.2cm}}
\toprule
№ & Замена & Ответ\\
\midrule
1 & $u=y-x,\ v=y$, $|J|=1$ & $14a^4$\\
2 & $x=a(t-\sin t),\ y=as(1-\cos t)$, \newline $|J|=a^2(1-\cos t)^2$ & $\dfrac{35\pi a^4}{12}$\\
3 & полярные, $|J|=r$ & $\int_{\pi/4}^{\pi/3}\!d\varphi\!\int_0^{2/\cos\varphi}\!f(r)r\,dr$ (второй порядок --- в задаче)\\
4 & полярные, $|J|=r$ & $-6\pi^2$\\
5 & $u=xy,\ v=y/x$, $|J|=\frac1{2v}$ & однократный интеграл (см.\ задачу) $=\dfrac{165}{128}-\ln2$\\
6 & $x=u,\ y=uv,\ z=u^2vw$, $|J|=u^3v$ & $\dfrac1{364}$\\
7 & сферические, $|J|=r^2\cos\psi$ & $\dfrac1{48}$\\
8 & $u=\frac z{y^2},\ v=\frac zx,\ w=z$, $|J|=\frac{w^{3/2}}{2u^{3/2}v^2}$ & $\dfrac{2}{27}\dfrac{(\sqrt b-\sqrt a)(\beta^3-\alpha^3)}{\sqrt{ab}\,\alpha^3\beta^3}h^4\sqrt h$\\
9 & цилиндрические, $|J|=r$ & $\dfrac3{35}$\\
10 & $x=u(1-v),\ y=uv$, $|J|=u$ & $\dfrac7{24}$\\
11 & цилиндрические, $|J|=r$ & $\dfrac{\pi a^3}{6}$\\
12 & сферические, $|J|=r^2\cos\psi$ & $\pi a^3$\\
13а & обобщ. сферические, $n=4$ & $\dfrac{abc}{90}$\\
13б & обобщ. сферические, $n=6$ & $\dfrac{abc}{1680}$\\
13в & обобщ. сферические, $n=3$ & $\dfrac{4\pi abc}{35}$\\
\bottomrule
\end{tabular}
\end{center}

Все ответы дополнительно проверены независимым способом (другая система координат, геометрия или метод сечений) --- см.\ строки «Проверка» в конце каждой задачи.
